\section{Benchmarks and results: consolidation}
\label{sec:consolidation}

This part puts the measures of the three services side by side: one host, the configurations measured, and where each
part of the system saturates.

\subsection{The host}

Every valid measure ran on one machine, recorded by the \code{host.txt} of each run:

\begin{table}[H]
\centering\small
\begin{tabularx}{\linewidth}{@{}L{3.4cm}X@{}}
\toprule
CPU & Intel Core i7-9750H, 2.60~GHz: 6 cores, 12 logical CPUs \\
Memory and storage & 16~GB; NVMe SSD \\
Operating system & Ubuntu 24.04.5, Linux 7.0 \\
Docker & 28.3.3, sees the 12 CPUs and 16~GB \\
JVM & OpenJDK 21.0.12 for JMH and the drivers; the images run Java 17 \\
Power & Plugged in, balanced energy profile \\
\bottomrule
\end{tabularx}
\caption{Reference machine of all the benchmarks.}
\label{tab:host}
\end{table}

The deployment limits each container (\code{.env}): MongoDB 2 CPUs and 1~GB with a 0.25~GB cache; Indexing 2 CPUs and
1~GB; Search 2 CPUs per replica; Ingestion 1 CPU and 512~MB; each Hazelcast member 1 CPU and 768~MB; the Control and
nginx 0.5 CPUs. The load experiments of Search give nginx 2 CPUs, so that it does not hide the rest, and measure its
0.5 CPUs apart. The limits together exceed the 12 CPUs of the host only in the Search experiments with 4 replicas
(14 CPUs), where the replicas compete for the same cores.

\subsection{What was measured}

\begin{table}[H]
\centering\small
\begin{tabularx}{\linewidth}{@{}L{3.1cm}L{2.2cm}L{2.2cm}X@{}}
\toprule
\textbf{Experiment} & \textbf{Scope} & \textbf{Replicas} & \textbf{Threads and clients} \\
\midrule
Ingestion JMH & component, host & -- & 1, 2, 4, 8 workers \\
Indexing JMH & component, host & -- & 1, 2, 4, 8 parser threads \\
Search JMH & component, host & -- & 1, 2, 4, 8 threads \\
Indexing, threads & Indexing + MongoDB + Hazelcast, Docker & 1 Indexing & 1, 2, 4, 8 parser threads; 2,000 books \\
Indexing, scale & same & 1 Indexing & 4 threads; 200, 10,000, 30,000 books \\
Indexing, publication & same & 1 Indexing & 1, 10, 50 new books on 10,000 \\
Search, load & Search + MongoDB + nginx, Docker & 1, 2, 4 Search & 1 to 256 closed-loop clients; 250 Jetty threads
  (also 16 and 64); pool of 50 (also 10 and 100); nginx with 2 CPUs (also 0.5); index of 15, 197 and 10,000 books \\
Search with the pipeline & complete pipeline, Docker & 1 of each & 32 and 128 clients on 2,000 books while the
  pipeline processes new ones \\
\bottomrule
\end{tabularx}
\caption{The experiments, each labelled with what it measures. Every configuration was repeated 3 times, except the
30,000 books of Indexing (once). Client concurrency, replicas and threads are kept apart.}
\label{tab:experiments}
\end{table}

\subsection{Bottlenecks and scaling trends}

\begin{table}[H]
\centering\small
\begin{tabularx}{\linewidth}{@{}L{1.9cm}L{4.4cm}X@{}}
\toprule
\textbf{Service} & \textbf{Measured rate} & \textbf{What limits it} \\
\midrule
Ingestion & 473 $\mu$s per book with 1 worker, 155 with 8 (local source, JMH) &
  The coordination: claim and complete cost 1.6~ms per book with 1 thread and 0.36 with 8, more than the ingestion
  itself. With Project Gutenberg the download is far longer than both. \\
Indexing & 15 books/s on 2,000 books; 5.5 on 10,000; 2.0 on 30,000 &
  MongoDB: the cost of \code{\$addToSet} grows with each posting list, and one writer sends one stream of writes.
  Threads and replicas of Indexing do not change it (Section~\ref{sec:benchmarks-indexing}). \\
Search & About 7,500 req/s on 15 books whatever the replicas; on 10,000 books 1,036 req/s with 1 replica and
  2,612 with 4 &
  On a small index MongoDB is the limit: it reaches its 2 CPUs while the replicas of Search have CPU to spare. On
  10,000 books each query matches most of the index, the CPU of Search is the limit, and replicas help: 4 give 2.5
  times the peak of one, less than 4 because the host runs out of CPU (Section~\ref{sec:benchmarks}). \\
The whole system & with the pipeline running, Search 2,600 req/s and Indexing 0.91 books/s at the same time &
  MongoDB is shared: the writes of Indexing take it from 80 to 191~\% of its 2 CPUs. Indexing drops from 2.42 to
  0.91 books/s while Search is under load, and the p95 of Search with 128 clients rises from 109 to 160~ms. \\
\bottomrule
\end{tabularx}
\caption{Where each part of the system saturates on the reference machine.}
\label{tab:bottlenecks}
\end{table}

\textbf{Threads and replicas.} Threads help where the work is CPU inside one process: the workers of Ingestion and the
parsers of Indexing in JMH, and the replicas of Search when queries are long. They do not help where the limit is a shared
resource: the queue of Hazelcast for Ingestion with small books, MongoDB for Indexing, and MongoDB again for Search
on a small index. Indexing replicas are
standbys by design, so their effect is failover, not throughput.

\textbf{What was not measured.} The full pipeline with several Ingestion replicas and Project Gutenberg as the source
(the measures use the local source, so that the network does not dominate); 100,000 books, because indexing them at the
measured rate would take most of a day; and two machines, which is Stage~3.
